% ============================================================
%  BEAMER — DERIVADAS APLICADAS A LA ECONOMÍA
%  Matemática Aplicada a la Economía · Temas 1–4
%  Autor: Ing. Gabriel Astudillo
%  Compilar:  latexmk -xelatex beamer_derivadas_economia.tex
%  Nota: la definición de derivada se escribe con \Delta x.
% ============================================================
\documentclass[aspectratio=169]{beamer}
\usepackage{mathwizards-beamer}
\mwbeamer{Matemática Aplicada a la Economía}{Derivadas}{Cálculo diferencial}

\begin{document}

\begin{frame}[plain]
  \titlepage
\end{frame}

% ════════════════════════════════════════════════════════════
\section{Marco Teórico}
% ════════════════════════════════════════════════════════════

\begin{frame}{Incrementos y Tasa de Cambio}
  \begin{block}{Incremento de la variable y de la función}
    \[
      \Delta x = x_2 - x_1
      \qquad\Longrightarrow\qquad
      x_2 = x_1 + \Delta x
    \]
    \[
      \Delta y = f(x_2) - f(x_1) = f(x_1 + \Delta x) - f(x_1)
    \]
  \end{block}
  \vspace{4pt}
  \begin{mwextra}[Tasa de cambio media]
    \[
      T = \frac{\Delta y}{\Delta x} = \frac{f(x_1 + \Delta x) - f(x_1)}{\Delta x}
    \]
    Es la \textbf{velocidad media} de cambio de $y$ respecto a $x$.
  \end{mwextra}
\end{frame}

\begin{frame}{Modelos Matemáticos en Economía}
  \begin{block}{Funciones económicas básicas}
    \begin{itemize}
      \item \textbf{Ingreso:} $I(q) = P(q)\cdot q$
      \item \textbf{Costo total:} $C(q) = C_v + C_f$
      \item \textbf{Beneficio:} $B(q) = I(q) - C(q)$
    \end{itemize}
  \end{block}
  \vspace{4pt}
  \begin{columns}[T]
    \column{0.48\textwidth}
    \begin{mwpasos}[Equilibrio de mercado]
      $S(P) = D(P)$: punto de intersección de oferta y demanda.
    \end{mwpasos}
    \column{0.48\textwidth}
    \begin{mwextra}[Interés compuesto]
      $P = P_0(1+r)^{n}$, con $n$ períodos y tasa $r$.
    \end{mwextra}
  \end{columns}
\end{frame}

\begin{frame}{La Derivada: Definición e Interpretación}
  \begin{block}{Tasa instantánea (definición)}
    \[
      f'(x) \;=\; \lim_{\Delta x \to 0} \frac{\Delta y}{\Delta x}
      \;=\; \lim_{\Delta x \to 0} \frac{f(x + \Delta x) - f(x)}{\Delta x}
    \]
    Requiere que $f$ sea \textbf{continua} en el punto.
  \end{block}
  \vspace{4pt}
  \begin{mwextra}[Interpretación geométrica]
    Cuando $\Delta x \to 0$, la recta \textbf{secante} tiende a la \textbf{tangente}:
    \[
      f'(x_0) = m_T = \text{pendiente de la tangente en } x_0
    \]
    Notaciones: $f'(x)$, $y'$, $\dfrac{dy}{dx}$, $\left(\dfrac{d}{dx}\right)y$.
  \end{mwextra}
\end{frame}

\begin{frame}{Reglas de Derivación y Derivadas Marginales}
  \begin{block}{Reglas fundamentales}
    \begin{columns}[T]
      \column{0.5\textwidth}
      \begin{itemize}
        \item $f = k \;\Rightarrow\; f' = 0$
        \item $f = x^{n} \;\Rightarrow\; f' = n\,x^{n-1}$
        \item $f = u\,v \;\Rightarrow\; f' = u'v + uv'$
        \item $f = \dfrac{u}{v} \;\Rightarrow\; f' = \dfrac{v\,u' - u\,v'}{v^{2}}$
      \end{itemize}
      \column{0.46\textwidth}
      \begin{itemize}
        \item $f = u^{n} \;\Rightarrow\; f' = n\,u^{n-1}u'$
        \item $f = e^{u} \;\Rightarrow\; f' = e^{u}u'$
        \item $f = \ln u \;\Rightarrow\; f' = \dfrac{u'}{u}$
      \end{itemize}
    \end{columns}
  \end{block}
  \vspace{2pt}
  \begin{mwpasos}[Derivadas marginales en economía]
    $I'(q)$ \textbf{ingreso marginal} \quad$\cdot$\quad $C'(q)$ \textbf{costo marginal} \quad$\cdot$\quad $B'(q)$ \textbf{beneficio marginal}
  \end{mwpasos}
\end{frame}

\begin{frame}{Técnicas Avanzadas}
  \begin{mwpasos}[Derivación logarítmica: $y = u(x)^{v(x)}$]
    \begin{enumerate}
      \item Aplicar $\ln$: $\ln y = v\,\ln u$ \quad (bajar exponente).
      \item Derivar: $\dfrac{y'}{y} = v'\,\ln u + v\,\dfrac{u'}{u}$.
      \item Despejar $y' = y\,[\;\cdots\;]$.
    \end{enumerate}
  \end{mwpasos}
  \vspace{2pt}
  \begin{columns}[T]
    \column{0.48\textwidth}
    \begin{mwextra}[Implícita]
      Derivar término a término con $y$ función de $x$ (aparece $y'$), agrupar y despejar.
    \end{mwextra}
    \column{0.48\textwidth}
    \begin{mwextra}[Sucesivas]
      $y' \;\to\; y'' \;\to\; y''' \;\to\; \dfrac{d^{n}y}{dx^{n}}$
    \end{mwextra}
  \end{columns}
\end{frame}

% ════════════════════════════════════════════════════════════
\section{Ejercicios de Clase}
% ════════════════════════════════════════════════════════════

\begin{frame}{Ejercicio 1}
  \begin{mwejercicio}[Ejercicio 1 — Incrementos y tasa \quad \medio]
    Para $f(x) = x^{2} + 3x$ con $x_1 = 1$ y $x_2 = 1{,}04$, calcule $\Delta x$, $\Delta y$ y la tasa de cambio media $T$.
  \end{mwejercicio}
  \vspace{3.5cm}
\end{frame}

\begin{frame}{Ejercicio 2}
  \begin{mwejercicio}[Ejercicio 2 — Interés compuesto \quad \medio]
    Se invierten $\$500$ a interés compuesto del $4\%$ anual. ¿Cuál es el capital al cabo de $5$ años?
  \end{mwejercicio}
  \vspace{3.5cm}
\end{frame}

\begin{frame}{Ejercicio 3}
  \begin{mwejercicio}[Ejercicio 3 — Derivada por definición ($\Delta x$) \quad \medio]
    Calcule $f'(x)$ por definición, con $\Delta x$, para $f(x) = x^{2} - 5x$.
  \end{mwejercicio}
  \vspace{3.5cm}
\end{frame}

\begin{frame}{Ejercicio 4}
  \begin{mwejercicio}[Ejercicio 4 — Regla del producto \quad \dificil]
    Derive $f(x) = (x^{2}+1)(3x-2)$.
  \end{mwejercicio}
  \vspace{3.5cm}
\end{frame}

\begin{frame}{Ejercicio 5}
  \begin{mwejercicio}[Ejercicio 5 — Cadena con raíz \quad \dificil]
    Derive $f(x) = \sqrt{5x^{2} + 2x}$.
  \end{mwejercicio}
  \vspace{3.5cm}
\end{frame}

\begin{frame}{Ejercicio 6}
  \begin{mwejercicio}[Ejercicio 6 — Logaritmo \quad \dificil]
    Derive $f(x) = \ln(3x^{2} - x)$.
  \end{mwejercicio}
  \vspace{3.5cm}
\end{frame}

\begin{frame}{Ejercicio 7}
  \begin{mwejercicio}[Ejercicio 7 — Derivación implícita \quad \dificil]
    Halle $y'$ si $x^{2} + 4xy - y^{2} = 0$.
  \end{mwejercicio}
  \vspace{3.5cm}
\end{frame}

\begin{frame}{Ejercicio 8}
  \begin{mwejercicio}[Ejercicio 8 — Derivación logarítmica \quad \muyDificil]
    Halle $y'$ por derivación logarítmica: $y = x^{x}$.
  \end{mwejercicio}
  \vspace{3.5cm}
\end{frame}

\begin{frame}{Ejercicio 9}
  \begin{mwejercicio}[Ejercicio 9 — Implícita con logaritmo \quad \muyDificil]
    Halle $y'$ si $3x^{2} + 2y^{3} - \ln y = 0$.
  \end{mwejercicio}
  \vspace{3.5cm}
\end{frame}

\begin{frame}{Ejercicio 10}
  \begin{mwejercicio}[Ejercicio 10 — Derivadas sucesivas \quad \muyDificil]
    Dada $y = A\,e^{3x} + B\,e^{x}$, demuestre que $y'' - 4y' + 3y = 0$.
  \end{mwejercicio}
  \vspace{3.5cm}
\end{frame}

\end{document}
